Let $\displaystyle{\alpha = \sqrt[3]{413 \mspace{1mu}} \mspace{2mu} - \mspace{-1mu} \sqrt[3]{3 \mspace{1mu}} \mspace{1mu}}$.
Then:
$\displaystyle{ \mspace{2mu}
  \alpha^3 \mspace{2mu}
  = \mspace{2mu} 413 - 3 \cdot \mspace{-2mu} \sqrt[3]{3} \cdot \mspace{-2mu} \sqrt[3]{413} \cdot \mspace{-2mu} \paren{\mspace{-2mu} \sqrt[3]{413} \mspace{2mu} - \mspace{-1mu} \sqrt[3]{3 \mspace{1mu}} \mspace{1mu}} - 3 \mspace{2mu}
  = \mspace{2mu} 410 - 3 \sqrt[3]{1239 \mspace{1mu}} \mspace{1mu} \alpha
\mspace{1mu}}$.

Therefore, $\alpha$ is a root of the polynomial function
$p\paren{x} = x^3 + 3 \sqrt[3]{1239 \mspace{1mu}} \mspace{1mu} x - 410$.
This function is strictly increasing, as a function $\mathbb{R} \to \mathbb{R}$, hence $\alpha$ is its only real root,
and, moreover, for every $x \in \mathbb{R}$,
$x < \alpha$ if and only $p\paren{x} < 0$,
and $x > \alpha$ if and only $p\paren{x} > 0$.

Putting $x = 6$, it needs to be determined whether $216 + 18 \sqrt[3]{1239} - 410$ is positive or negative,
which simplifies to determining whether $97$ is, respectively, less or greater than $9 \sqrt[3]{1239}$.
Now:
\begin{align*}
  \paren{\frac{97}{9}}^3
  &= \paren{10 + \frac{7}{9}}^3
   > 1000 + 3 \cdot 100 \cdot \frac{7}{9} + 3 \cdot 10 \cdot \frac{49}{81} = \\
  &= 1000 + \frac{700}{3} + \frac{490}{27}
   > 1000 + 233 + 10 > \\
  &> 1239
\end{align*}
It follows that $97 > 9 \sqrt[3]{1239 \mspace{1mu}} \mspace{1mu}$, so $p\paren{6} < 0$,
which means that $6 < \alpha = \sqrt[3]{413} \mspace{2mu} - \mspace{-1mu} \sqrt[3]{3}$.
So, in conclusion:
\begin{equation*}
  6 \mspace{2mu} + \mspace{-1mu} \sqrt[3]{3 \mspace{1mu}} \mspace{2mu}
  <
  \sqrt[3]{413 \mspace{1mu}}
\end{equation*}
